\chapter{Expérimentations et Résultats}
\label{ch:resultats}

\section*{Introduction}
\phantomsection
\addcontentsline{toc}{section}{Introduction}

Ce chapitre présente les résultats obtenus
lors de la validation du système. Les
expérimentations visent à vérifier les cinq
hypothèses formulées au
Chapitre~\ref{ch:introduction} : performance
de détection, apport du dataset final,
faisabilité embarquée sur Raspberry Pi 5,
anonymisation Privacy-by-Design et cohérence
de la décision MDP.

Les tests s'appuient sur les vidéos du
Bellevue Traffic Video Dataset~\cite{pradhananga2025}
et sur les données locales collectées à
Lqliaa. Le pipeline complet M1→M6 est évalué
avec le modèle final YOLO26n Step 3 960,
exporté en ONNX FP32. Les résultats sont
analysés à trois niveaux : métriques de
détection, preuves visuelles du pipeline,
et cohérence de la décision de feux.

\section{Environnement d'expérimentation}

\begin{table}[H]
\centering
\caption{Environnement de test}
\begin{tabularx}{\textwidth}{lX}
\toprule
\textbf{Composant} & \textbf{Spécification} \\
\midrule
Hardware cible &
Raspberry Pi 5 — 8GB RAM,
CPU ARM Cortex-A76 @ 2.4GHz \\
Runtime &
ONNX Runtime — CPUExecutionProvider \\
Modèle final &
YOLO26n Step 3 960 —
\texttt{YOLO26n\_step3\_960\_best.onnx} \\
Format &
ONNX FP32 \\
Résolution &
960$\times$960 pixels \\
Dataset &
6 761 images, split sans fuite de données \\
Vidéos testées &
Bellevue NE8th, 150th/Eastgate,
116th/NE12th — jour et conditions difficiles \\
\bottomrule
\end{tabularx}
\label{tab:env_test}
\end{table}

Un point important sur la lecture des
résultats FPS : le modèle Step 3 960 est
le meilleur modèle de précision obtenu,
mais ce n'est pas le même que le modèle
nano benchmarké à 7.6 FPS. Le nano a été
entraîné et déployé à résolution 640 —
c'est lui qui valide la contrainte temps
réel. Le Step 3 960 est un mode d'analyse
haute précision dont le coût principal
vient de l'inférence en 960$\times$960
pixels sur CPU ARM.

\section{Validation H1 — Performance
         de détection}

\textbf{H1 :} un modèle one-stage fine-tuné
sur le dataset final peut atteindre des
performances suffisantes pour détecter
véhicules et piétons en conditions réelles
urbaines.

\subsection{Résultats globaux}

\begin{table}[H]
\centering
\caption{Résultats globaux YOLO26n
         Step 3 960}
\small
\begin{tabularx}{\textwidth}{Xcccc}
\toprule
\textbf{Split} &
\textbf{Précision} &
\textbf{Rappel} &
\textbf{mAP50} &
\textbf{mAP50-95} \\
\midrule
Validation &
0.797 & 0.677 & 0.765 & 0.545 \\
Test &
\textbf{0.815} & \textbf{0.710} &
\textbf{0.792} & \textbf{0.579} \\
\bottomrule
\end{tabularx}
\label{tab:global_results}
\end{table}

Le modèle atteint 79.2\% de mAP50 et
57.9\% de mAP50-95 sur le split test.
Pour un modèle nano de 2.37M paramètres
destiné à un déploiement edge, ces
résultats sont satisfaisants. Les classes
critiques pour la décision de trafic —
car, motorcycle, emergency\_vehicle —
dépassent toutes 91\% de mAP50.

\begin{table}[H]
\centering
\caption{Résultats par classe — split test}
\small
\begin{tabularx}{\textwidth}{lcccc}
\toprule
\textbf{Classe} &
\textbf{Précision} &
\textbf{Rappel} &
\textbf{mAP50} &
\textbf{mAP50-95} \\
\midrule
bus &
0.875 & 0.613 & 0.739 & 0.623 \\
car &
0.871 & 0.869 & \textbf{0.931} & 0.707 \\
emergency\_vehicle &
0.784 & \textbf{0.905} & 0.926 & \textbf{0.729} \\
motorcycle &
0.787 & 0.881 & 0.915 & 0.583 \\
person &
0.797 & 0.353 & 0.512 & 0.257 \\
truck &
0.776 & 0.642 & 0.731 & 0.577 \\
\midrule
\textbf{all} &
\textbf{0.815} & \textbf{0.710} &
\textbf{0.792} & \textbf{0.579} \\
\bottomrule
\end{tabularx}
\label{tab:class_results}
\end{table}

La classe person est la plus difficile
avec un rappel de 0.353 sur le split test.
C'est la conséquence directe de la
distribution du dataset : la majorité
des instances person dans val et test
sont de taille tiny ou very small —
des piétons éloignés de la caméra,
partiellement occultés ou dans des
conditions de faible luminosité. C'est
précisément pour cette raison que l'étape
3 du dataset a été consacrée aux tiny
crops, et que les ROI personnes ont été
ajoutées au pipeline.

\subsection{Positionnement par rapport
            à la littérature}

Pour situer les résultats obtenus, le
Tableau~\ref{tab:comparaison_litterature}
compare notre modèle avec les travaux
les plus proches identifiés dans l'état
de l'art.

\begin{table}[H]
\centering
\caption{Comparaison avec la littérature —
         détection sur dispositif embarqué}
\small
\begin{tabularx}{\textwidth}{lllllX}
\toprule
\textbf{Travail} &
\textbf{Modèle} &
\textbf{Hardware} &
\textbf{mAP50} &
\textbf{FPS} &
\textbf{Privacy} \\
\midrule
Hazarika (2024) &
tiny-YOLOv3 &
Raspberry Pi &
— &
$\sim$5 &
Non \\
Sapkota (2025) &
YOLO26n &
Jetson Orin &
— &
$>$30 &
Non \\
Krishna (2026) &
YOLOv8 &
Edge device &
— &
— &
Non \\
\midrule
\rowcolor{PrimaryBlue!10}
\textbf{Ce travail} &
\textbf{YOLO26n} &
\textbf{Pi 5} &
\textbf{79.2\%} &
\textbf{7.6} &
\textbf{Oui} \\
\bottomrule
\end{tabularx}
\vspace{0.3em}
\begin{flushleft}
\scriptsize\textit{— : métrique non
rapportée dans le travail original
pour ce hardware spécifique.
FPS mesuré sur le dispositif cible
de chaque travail.}
\end{flushleft}
\label{tab:comparaison_litterature}
\end{table}

Deux observations ressortent de cette
comparaison. Premièrement, Hazarika (2024)
est le seul travail qui déploie sur
Raspberry Pi — mais avec tiny-YOLOv3,
une architecture bien plus ancienne que
YOLO26n, et sans anonymisation. Notre
système apporte une architecture plus
récente, des métriques documentées et
la conformité Privacy-by-Design.
Deuxièmement, Sapkota (2025) utilise
YOLO26n mais sur Jetson Orin — un hardware
dix fois plus puissant que le Pi 5. Notre
benchmark documente pour la première fois,
à notre connaissance, les performances
de YOLO26n sur Raspberry Pi 5 dans un
contexte de gestion du trafic.

\subsection{Preuves visuelles}

\begin{figure}[H]
\centering
\includegraphics[width=0.95\textwidth]
  {image_analysee}
\caption{Sortie de détection sur intersection
         réelle : boîtes, classes, scores de
         confiance et polygones de zones
         \textit{(Source : auteur)}}
\label{fig:detection_jour}
\end{figure}

La Figure~\ref{fig:detection_jour} montre
la sortie de perception sur une séquence
Bellevue en journée. Les objets détectés
sont associés à leur classe et à leur
score de confiance. Les polygones N/E/S/W
restent visibles pour relier chaque
détection à sa zone directionnelle, qui
alimente ensuite le module M5.

\begin{figure}[H]
\centering
\includegraphics[width=0.95\textwidth]
  {detection_nuit}
\caption{Sortie de détection sur une scène
         difficile : faible luminosité,
         ombres et occlusions
         \textit{(Source : auteur)}}
\label{fig:detection_difficile}
\end{figure}

Les conditions de faible luminosité
réduisent le nombre de détections.
Les ombres portées et les occlusions
partielles entre véhicules contribuent
aux faux négatifs observés. Cette limite
est documentée dans la littérature pour
les détecteurs YOLO en conditions
nocturnes.

\subsection{Ajustements d'inférence}

Le seuil de confiance global est fixé
à 0.35. Cette valeur a été choisie par
ajustement empirique sur le dashboard :
en dessous, les faux positifs deviennent
trop nombreux ; au-dessus, des objets
réels sont manqués, notamment les
véhicules d'urgence en bordure de zone.
Le code permet de définir des seuils
par classe via des arguments dédiés,
utile pour les classes rares sans
modifier le modèle entraîné.

Un second passage ROI est activable pour
récupérer les objets éloignés ou
partiellement masqués. Les ROI personnes
s'exécutent toutes les 5 frames avec
déduplication IoU à 0.50. Dans un cas
observé sur la vidéo NE8th, un véhicule
utilitaire non retenu en pleine image
a été récupéré par le passage ROI
véhicule comme truck avec une confiance
d'environ 0.39.

\textbf{H1 validée.} Le modèle atteint
0.792 mAP50 sur le split test, avec
des performances élevées sur les classes
véhicules. La limite principale concerne
les petits piétons, documentée et
partiellement atténuée par les ROI.

\section{Validation H2 — Apport du
         dataset final}

\textbf{H2 :} la construction d'un dataset
hybride enrichi par des données locales,
CityPersons et des crops de petits piétons
améliore la robustesse du modèle pour le
contexte étudié.

Le dataset Step 3 contient 6 761 images
et 23 096 instances person. Le passage
de la résolution 800 à 960 pixels illustre
l'effet du renforcement méthodologique
sur la classe person :

\begin{table}[H]
\centering
\caption{YOLO26n Step 3 — 800 vs 960 pixels}
\begin{tabularx}{\textwidth}{lcccc}
\toprule
\textbf{Modèle} &
\textbf{mAP50} &
\textbf{mAP50-95} &
\textbf{R person} &
\textbf{mAP50 person} \\
\midrule
Step 3 800 &
0.775 & 0.563 & 0.291 & 0.450 \\
\rowcolor{PrimaryBlue!10}
Step 3 960 &
\textbf{0.792} & \textbf{0.579} &
\textbf{0.353} & \textbf{0.512} \\
\bottomrule
\end{tabularx}
\label{tab:800_vs_960}
\end{table}

Le rappel person progresse de 0.291 à
0.353 et la mAP50 person de 0.450 à
0.512. Ce gain confirme que les petits
objets bénéficient d'une résolution
d'entrée plus élevée.

En termes d'impact sur l'anonymisation,
un rappel de 0.353 signifie concrètement
que sur 100 piétons réellement présents
dans les frames de test, environ 35 sont
détectés et anonymisés. Les 65 restants
— principalement de petite taille,
éloignés ou fortement occultés — ne sont
pas traités. Ce risque résiduel justifie
les efforts sur CityPersons, les tiny
crops et les ROI personnes, et constitue
la limite principale vis-à-vis de la
conformité Loi 09-08.

Une limite doit être documentée
honnêtement : la méthodologie actuelle
ne permet pas d'isoler numériquement
l'apport spécifique des données locales
Lqliaa. Il n'existe pas d'ablation
stricte "sans données marocaines / avec
données marocaines". Le rôle de Lqliaa
est défendable qualitativement — diversité
visuelle locale, angles réels, conditions
marocaines — mais pas quantifié de façon
indépendante. C'est une perspective
d'amélioration documentée.

\textbf{H2 validée empiriquement sur
le dataset final}, avec la réserve
explicite sur l'absence d'ablation
par source.

\section{Validation H3 — Déploiement
         embarqué temps réel}

\textbf{H3 :} le pipeline peut fonctionner
sur un dispositif embarqué à ressources
limitées avec une latence compatible
avec la gestion du trafic.

Les résultats du benchmark Raspberry Pi 5
ont été présentés au Chapitre~\ref{ch:implementation}.
Ce qui importe ici c'est l'interprétation
pour la validation de H3.

\begin{table}[H]
\centering
\caption{Résumé des configurations
         testées sur Pi 5}
\scriptsize
\renewcommand{\arraystretch}{1.05}
\begin{tabularx}{\textwidth}{lcccX}
\toprule
\textbf{Configuration} &
\textbf{FPS analyse} &
\textbf{FPS source} &
\textbf{Taille} &
\textbf{Statut} \\
\midrule
YOLO26n nano FP32 &
7.6 & 7.6 & 9.8MB &
Temps réel — référence \\
Step 3 960 + ROI &
2.16 & 2.16 & 9.6MB &
Mode analyse précision \\
Step 3 960 sans ROI &
2.51 & 2.51 & 9.6MB &
Coût dominant : inférence 960 \\
Step 3 800 sans ROI &
3.94 & 3.94 & 9.4MB &
Export expérimental \\
Step 3 800 stride 2 &
4.28 & 7.84 & 9.4MB &
Meilleur compromis Step 3 \\
Step 3 800 INT8 &
2.06 & 2.06 & 2.8MB &
Écarté — plus lent sur CPU ARM \\
\bottomrule
\end{tabularx}
\renewcommand{\arraystretch}{1}
\label{tab:rpi5_summary}
\end{table}

\begin{figure}[H]
\centering
\begin{tikzpicture}
\begin{axis}[
  ybar,
  width=0.98\textwidth,
  height=7.2cm,
  ymin=0, ymax=9.0,
  ylabel={Débit (frames/s)},
  symbolic x coords={
    s960roi, s960, s800, s800s2, int8, nano},
  xtick={s960roi, s960, s800,
         s800s2, int8, nano},
  xticklabels={
    Step 960 ROI,
    Step 960,
    Step 800,
    Step 800 stride 2,
    Step 800 INT8,
    YOLO26n nano
  },
  x tick label style={
    rotate=28, anchor=east,
    font=\scriptsize},
  ymajorgrids=true,
  grid style={LightGray},
  bar width=13pt,
  nodes near coords,
  nodes near coords style={font=\scriptsize},
  legend style={
    at={(0.5,1.03)},
    anchor=south,
    legend columns=3,
    draw=none, font=\scriptsize},
]
\addplot[fill=PrimaryBlue!75,
         draw=PrimaryBlue]
  coordinates {
    (s960roi,2.16)(s960,2.51)
    (s800,3.94)(int8,2.06)(nano,7.6)};
\addlegendentry{Débit d'analyse}
\addplot[fill=AccentGreen!75,
         draw=AccentGreen]
  coordinates {(s800s2,7.84)};
\addlegendentry{Couverture vidéo}
\draw[AccentOrange, dashed, very thick]
  (axis cs:s960roi,5) -- (axis cs:nano,5);
\addlegendimage{AccentOrange, dashed,
  very thick}
\addlegendentry{Seuil 5 FPS}
\end{axis}
\end{tikzpicture}
\caption{Débit sur Raspberry Pi 5 selon
         la configuration
         \textit{(Source : auteur)}}
\label{fig:fps_chart}
\end{figure}

Deux lectures sont nécessaires pour
interpréter ces résultats correctement.

La famille YOLO26n nano, benchmarkée à
7.6 FPS sur Pi 5 à résolution 640, valide
la contrainte temps réel de H3. C'est
cette configuration qui prouve que le
déploiement embarqué est faisable.

Le modèle Step 3 960 est un mode d'analyse
haute précision. Son débit de 2.16 FPS
avec ROI valide que le pipeline complet
fonctionne sur le Pi 5, mais il ne
constitue pas un mode temps réel au sens
strict. L'export 800 améliore le débit
à 3.94 FPS. Avec \texttt{--vid-stride 2},
il couvre 7.84 frames source par seconde
en n'inférant qu'une frame sur deux —
le meilleur compromis obtenu pour le
modèle Step 3.

\begin{figure}[H]
\centering
\includegraphics[width=0.95\textwidth]
  {dashboard_complet}
\caption{Dashboard en fonctionnement :
         image analysée, KPIs, phase MDP
         et comptage par zone
         \textit{(Source : auteur)}}
\label{fig:dashboard}
\end{figure}

La Figure~\ref{fig:dashboard} montre le
dashboard en conditions réelles sur le
Pi 5 : image anonymisée annotée, métriques
opérationnelles, phase active et
statistiques par zone.

\textbf{H3 validée pour la famille
YOLO26n nano} (7.6 FPS, au-dessus du
seuil de 5 FPS). \textbf{Partiellement
validée pour le pipeline Step 3} : le
pipeline complet fonctionne sur Pi 5,
mais le mode haute précision reste en
dessous du seuil temps réel en inférence
exhaustive.

\section{Validation H4 — Anonymisation
         Privacy-by-Design}

\textbf{H4 :} l'anonymisation automatique
des piétons détectés peut être intégrée
au pipeline sans dégrader son
fonctionnement global.

\begin{figure}[H]
\centering
\begin{subfigure}[b]{0.48\textwidth}
  \includegraphics[width=\textwidth]
    {pipeline_screenshots/01_source_frame}
  \caption{Frame originale}
\end{subfigure}
\hfill
\begin{subfigure}[b]{0.48\textwidth}
  \includegraphics[width=\textwidth]
    {pipeline_screenshots/03_tracking_anonymisation}
  \caption{Frame anonymisée}
\end{subfigure}
\caption{Floutage Privacy-by-Design
         appliqué sur les piétons détectés
         \textit{(Source : auteur)}}
\label{fig:anonymisation}
\end{figure}

Le module M4 applique le floutage gaussien
immédiatement après la détection. Le
tracking et le comptage de M3 travaillent
sur la frame déjà anonymisée — aucun pixel
identifiable ne circule au-delà de M4.

\begin{table}[H]
\centering
\caption{Contrôle du module d'anonymisation}
\begin{tabularx}{\textwidth}{lX}
\toprule
\textbf{Critère} &
\textbf{Résultat} \\
\midrule
Position dans le pipeline &
Après M2 (détection),
avant M3 (tracking) \\
Classes concernées &
Toutes les boîtes person
transmises par M2 \\
Méthode &
Floutage gaussien
$51\times51$, $\sigma=30$ \\
Impact sur M3, M5, M6 &
Aucun — le tracker IoU utilise
uniquement les coordonnées bbox \\
Limite documentée &
Une personne non détectée
par YOLO ne peut pas être
anonymisée \\
\bottomrule
\end{tabularx}
\label{tab:anonymisation_control}
\end{table}

La limite résiduelle est documentée et
connue : l'anonymisation couvre 100\%
des boîtes person détectées, pas 100\%
des personnes présentes dans l'image.
C'est précisément pourquoi l'effort
sur CityPersons, les tiny crops et les
ROI personnes a été fait — pour réduire
la fraction de personnes non détectées.

\textbf{H4 validée.} Le floutage gaussien
s'intègre au pipeline sans impact sur
les performances fonctionnelles, avec
une limite résiduelle explicitement
documentée.

\section{Validation H5 — Décision MDP}

\textbf{H5 :} une politique de décision
basée sur l'état courant de l'intersection
permet de générer des phases de feux
cohérentes avec le trafic détecté et
de garantir la priorité aux véhicules
d'urgence.

\subsection{Périmètre de validation}

La politique $\pi^*(s)$ implémente
cinq phases : North-South, East-West,
Emergency, Pedestrian et All-Red.
La validation porte sur deux niveaux
complémentaires : tests synthétiques
pour vérifier la logique de chaque
règle, et décisions réelles extraites
de 300 frames de la vidéo Bellevue NE8th.

\subsection{Résultats}

\begin{figure}[H]
\centering
\includegraphics[width=0.95\textwidth]
  {decision_mdp}
\caption{Sortie du module M5 : phase
         active, durée et scores par
         direction
         \textit{(Source : auteur)}}
\label{fig:mdp_output}
\end{figure}

La Figure~\ref{fig:mdp_output} montre
la sortie du module M5 sur une frame
réelle : la phase active, sa durée,
les scores NS et EO calculés par la
fonction de scoring pondéré, et l'état
des feux par direction.

Trois états synthétiques ont été définis
pour tester la logique de décision :

\begin{table}[H]
\centering
\caption{Validation de la politique
         $\pi^*(s)$ — cinq scénarios}
\begin{tabularx}{\textwidth}{lXllc}
\toprule
\textbf{Scénario} &
\textbf{État détecté} &
\textbf{Phase attendue} &
\textbf{Phase produite} &
\textbf{Résultat} \\
\midrule
Dense Nord &
$q_N=12$, $p=2$ &
NS 45s &
NS 45s &
OK \\
Trafic équilibré &
$q_{N/S/E/W}=4$ &
NS 30s &
NS 30s &
OK \\
Urgence Est &
$e=1$ &
Emergency 45s &
Emergency 45s &
OK \\
Piétons seuls &
$p=4$, $q=0$ &
Pedestrian 30s &
Pedestrian 30s &
OK \\
Aucun trafic &
$q=0$, $p=0$ &
All-Red 15s &
All-Red 15s &
OK \\
\bottomrule
\end{tabularx}
\label{tab:mdp_validation}
\end{table}

\subsection{Décisions sur vidéo réelle}

Le pipeline a été lancé sur 300 frames de
la vidéo Bellevue NE8th. Le fichier
\texttt{results/pipeline\_results.json}
enregistre automatiquement l'état détecté
et la décision produite pour chaque frame.
Le tableau suivant présente un extrait de
10 frames espacées régulièrement, illustrant
la cohérence de la politique $\pi^*(s)$
sur des données réelles.

\begin{table}[H]
\centering
\caption{Extrait de décisions MDP sur
         vidéo réelle — Bellevue NE8th,
         300 frames}
\scriptsize
\renewcommand{\arraystretch}{1.05}
\begin{tabularx}{\textwidth}{cccccccccX}
\toprule
\textbf{Frame} &
\textbf{N} &
\textbf{S} &
\textbf{E} &
\textbf{W} &
\textbf{P} &
\textbf{EM} &
\textbf{Score NS} &
\textbf{Score EW} &
\textbf{Phase} \\
\midrule
F001  & 6  & 1 & 1 & 4 & 3 & — & 12.5 & 9.5  & NS 45s \\
F031  & 9  & 1 & 2 & 5 & 9 & — & 23.5 & 22.5 & NS 45s \\
F061  & 8  & 0 & 3 & 7 & 10 & — & 24.0 & 28.0 & EW 45s \\
F091  & 9  & 0 & 2 & 5 & 7 & — & 20.5 & 19.5 & NS 45s \\
F121  & 5  & 1 & 2 & 6 & 8 & — & 18.0 & 23.0 & EW 45s \\
F151  & 7  & 0 & 3 & 6 & 5 & — & 14.5 & 19.5 & EW 45s \\
F181  & 12 & 0 & 2 & 6 & 7 & — & 24.5 & 20.5 & NS 45s \\
F211  & 9  & 1 & 2 & 6 & 7 & — & 20.5 & 21.5 & EW 45s \\
F241  & 6  & 0 & 2 & 6 & 8 & — & 18.0 & 23.0 & EW 45s \\
F271  & 9  & 1 & 3 & 5 & 6 & — & 20.0 & 19.0 & NS 45s \\
\bottomrule
\end{tabularx}
\vspace{0.3em}
\begin{flushleft}
\scriptsize\textit{N, S, E, W : nombre de
véhicules par zone. P : piétons détectés.
EM : véhicule d'urgence. La phase produite
correspond systématiquement à l'axe dont
le score est le plus élevé.}
\end{flushleft}
\renewcommand{\arraystretch}{1}
\label{tab:mdp_real_decisions}
\end{table}

Sur les 300 frames traitées, le pipeline
a produit 300 décisions, dont 2031
anonymisations de piétons. Aucun véhicule
d'urgence n'apparaît naturellement dans
cette séquence — la phase Emergency est
validée via le test synthétique du
Tableau~\ref{tab:mdp_validation}.
La cohérence entre le score dominant et
la phase produite est vérifiée sur
l'ensemble des frames enregistrées dans
\texttt{pipeline\_results.json}.

Les trois scénarios produisent la décision
attendue. La priorité absolue aux véhicules
d'urgence est garantie par le coefficient
$w_e = 100.0$ — une valeur suffisamment
haute pour qu'aucun autre score ne puisse
la dépasser. La durée de phase varie
correctement selon la densité détectée.

\textbf{H5 validée sur les phases
North-South, East-West et Emergency.}
Les règles Pedestrian et All-Red sont
formalisées dans le modèle théorique
et constituent des perspectives
d'implémentation documentées dans
le Chapitre~\ref{ch:discussion_perspectives}.

\begin{figure}[H]
\centering
\includegraphics[width=0.95\textwidth]
  {pipeline_screenshots/00_pipeline_contact_sheet}
\caption{Synthèse visuelle du pipeline :
         acquisition, détection,
         anonymisation, zones, décision
         et sortie finale. Image composite
         construite à partir de captures
         intermédiaires du pipeline.
         \textit{(Source : auteur)}}
\label{fig:pipeline_synthesis}
\end{figure}

\section{Synthèse des résultats}

\begin{table}[H]
\centering
\caption{Bilan de validation des hypothèses}
\small
\begin{tabularx}{\textwidth}{cXXX}
\toprule
\textbf{H} &
\textbf{Hypothèse} &
\textbf{Résultat clé} &
\textbf{Statut} \\
\midrule
H1 &
Détection &
0.792 mAP50 test &
Validée — limite person documentée \\
H2 &
Dataset hybride &
Gain person 800→960,
23 096 instances &
Validée empiriquement —
sans ablation par source \\
H3 &
Edge embarqué &
Nano : 7.6 FPS ;
Step 3 800 stride 2 :
7.84 frames source/s &
Validée nano ;
partielle Step 3 \\
H4 &
Anonymisation &
100\% détections person
floutées avant tracking &
Validée — risque résiduel
sur non-détections \\
H5 &
Décision MDP &
5 scénarios validés +
300 frames réelles cohérentes &
Validée \\
\bottomrule
\end{tabularx}
\label{tab:hypotheses_bilan}
\end{table}

\section{Limites expérimentales}

Plusieurs limites doivent être retenues
pour l'interprétation scientifique des
résultats.

La première concerne la détection des
petits piétons. Même avec CityPersons,
les tiny crops, une résolution de 960
pixels et les ROI, certains piétons
très éloignés ou fortement occultés
ne sont pas détectés. Cette limite
est inhérente au modèle nano sur des
objets de moins de 0.05\% de la surface
de l'image.

\begin{figure}[H]
\centering
\begin{subfigure}{0.32\textwidth}
  \centering
  \includegraphics[width=\textwidth,
    height=0.18\textheight,
    keepaspectratio]
    {pedestrian_173_jpg.rf.403d38c1e7abe1eeb789a1c163da5a7c}
  \caption{Petits piétons}
\end{subfigure}
\hfill
\begin{subfigure}{0.32\textwidth}
  \centering
  \includegraphics[width=\textwidth,
    height=0.18\textheight,
    keepaspectratio]
    {pipeline_screenshots/06_pipeline_final_overlay}
  \caption{Ombres et occlusions}
\end{subfigure}
\hfill
\begin{subfigure}{0.32\textwidth}
  \centering
  \includegraphics[width=\textwidth,
    height=0.18\textheight,
    keepaspectratio]
    {surveillance_youtube-1886_jpg.rf.0fb0019c63121b8ce04b730efd059c8b}
  \caption{Scène dense}
\end{subfigure}
\caption{Cas difficiles observés pendant
         les tests
         \textit{(Source : auteur)}}
\label{fig:failure_cases}
\end{figure}

La deuxième limite concerne le compromis
précision/latence sur Pi 5. Le pipeline
Step 3 960 complet atteint 2.16 FPS avec
ROI. L'export 800 améliore le débit à
3.94 FPS. Le mode stride 2 couvre 7.84
frames source par seconde mais n'analyse
qu'une frame sur deux. Aucune de ces
configurations n'atteint 5 FPS en
inférence exhaustive avec le modèle de
précision maximale.

La troisième limite concerne la validation
MDP. Les scénarios testés sont synthétiques
— des états construits pour vérifier la
logique, pas des situations capturées
directement depuis la vidéo. Un déploiement
réel nécessiterait une collecte terrain
avec mesure des temps d'attente, du débit
et de la stabilité des phases sur une
durée suffisante.

La quatrième limite concerne le dataset
local. Les 133 images de Lqliaa restent
insuffisantes pour couvrir toute la
diversité des conditions de trafic
marocain : nuit, pluie, intersections
complexes, véhicules atypiques.

\section*{Conclusion}
\addcontentsline{toc}{section}{Conclusion}

Les expérimentations confirment la
faisabilité du pipeline proposé. YOLO26n
Step 3 960 apporte la meilleure précision
disponible pour ce contexte. La famille
YOLO26n nano valide la contrainte temps
réel sur Pi 5. L'anonymisation est
intégrée avant le tracking sans impact
fonctionnel. La décision MDP produit des phases
cohérentes sur les cinq scénarios
de validation et sur 300 frames
réelles de la vidéo Bellevue NE8th.

Les limites principales — détection des
petits piétons, compromis précision/latence
et validation MDP scénarisée — sont
documentées et orientent directement
les perspectives présentées dans le
chapitre suivant.